\documentclass[12pt]{article}
\usepackage[T1]{fontenc}
\usepackage[utf8]{inputenc}
\usepackage{lmodern}
\usepackage{textcomp}
\usepackage{enumitem}
\usepackage{geometry}
\usepackage{graphicx}
\usepackage{amsmath, amssymb}
\geometry{margin=1in}
\begin{document}
\begin{center}
Astronomy - Graduate Level Assessment 2 \
Total Marks: 40 \
Answer all questions.
\end{center}

\section*{Multiple Choice (10 marks)}
\begin{enumerate}[label=\arabic*.]
\item Imagine that the Earth was instantly moved to an orbit three times further away from the Sun. How much longer would a year be?
\begin{enumerate}[label=(\alph*)]
    \item exactly 3 times longer
    \item about 5.2 times longer
    \item Not enough information. It will depend on the inclination of the new orbit
    \item The length of the year wouldn't change because the Earth's mass stays the same.
\end{enumerate}
\item Say the pupil of your eye has a diameter of 5 mm and you have a telescope with an aperture of 50 cm. How much more light can the telescope gather than your eye?
\begin{enumerate}[label=(\alph*)]
    \item 10000 times more
    \item 100 times more
    \item 1000 times more
    \item 10 times more
\end{enumerate}
\item Which of these has NOT been one of the main hypotheses considered for the origin of the Moon?
\begin{enumerate}[label=(\alph*)]
    \item The Moon split from the Earth due to tidal forces.
    \item The Moon was captured into Earth orbit.
    \item The Earth and Moon co-accreted in the solar nebula.
    \item Earth was rotating so rapidly that the Moon split from it.
\end{enumerate}
\item Which one of these constellations is not located along the Milky Way in the sky?
\begin{enumerate}[label=(\alph*)]
    \item Perseus
    \item Cygnus
    \item Scorpius
    \item Leo
\end{enumerate}
\item In addition to the conditions required for any solar eclipse what must also be true in order for you to observe a total solar eclipse?
\begin{enumerate}[label=(\alph*)]
    \item The Earth must lie completely within the Moon's penumbra.
    \item The Moon's penumbra must touch the area where you are located.
    \item The Earth must be near aphelion in its orbit of the Sun.
    \item The Moon's umbra must touch the area where you are located.
\end{enumerate}
\end{enumerate}

\section*{True / False (10 marks)}
\textit{Answer True or False}
\begin{enumerate}[label=\arabic*.]
\setcounter{enumi}{5}
\item True or False: Water-ice is distributed evenly across Mercury, rather than being confined to the polar regions.
\item True or False: The Andromeda Galaxy is located approximately 2.5 million light years away from Earth.
\item True or False: The force required to accelerate the truck remains constant regardless of the planetary location due to uniform mass.
\item True or False: The resolution of a telescope is a measure of its ability to measure the angular separation of closely spaced objects.
\item True or False: The Moon's surface is less impacted by craters due to its thicker atmosphere compared to Earth's atmosphere.
\end{enumerate}

\section*{Long Form Response (20 marks)}
\textit{Answer each question with a clear and complete explanation.}
\begin{enumerate}[label=\arabic*.]
\setcounter{enumi}{10}
\item Discuss the reasons behind the varying compositions of inner and outer planets in our solar system, focusing on the role of temperature and material condensation in the protoplanetary nebula.
\item Discuss the concept of bolometric luminosity in astronomy, explaining its significance and the methods used to integrate luminosity across all wavelengths.
\end{enumerate}

\vfill
\noindent\textit{End of Paper}
\end{document}